\documentclass[10pt]{article}
\usepackage[T1]{fontenc}
\usepackage{lmodern,amsmath,amssymb,amsthm,microtype}
\usepackage[margin=.8in]{geometry}
\usepackage[colorlinks=true,urlcolor=blue]{hyperref}
\newtheorem{theorem}{Theorem}
\newtheorem{lemma}{Lemma}
\title{An Euler Product over the Inert Primes of $\mathbb Q(\sqrt3)$}
\author{Qichao Wang\\{\small Hebei University of Technology}\\[.3em]
Daoyu Dong\\{\small University of Electronic Science and Technology of China}}
\date{October 27, 2026}
\begin{document}\maketitle
\begin{abstract}
We prove the prime-product conjecture in OEIS A003630. The partial
product is exactly $9/8$ times a quotient of two partial Dirichlet
$L$-products. Their convergence in increasing prime order, together
with two elementary $L$-value integrals, gives the conjectured limit
$\sqrt3/2$.
\end{abstract}

\section{The product and its character identity}
The inert rational primes of $\mathbb Q(\sqrt3)$ are exactly the
primes congruent to $5$ or $7$ modulo $12$. This classification is
already established in A003630~\cite{oeis}, including Jianing Song's 2018
comment. The separate conjecture of Dimitris Valianatos, dated
March 26, 2017, concerns the following product.

\begin{theorem}\label{thm:main}
With factors taken in increasing order of the primes,
\[
 \prod_{\substack{p\text{ prime}\\p\equiv5,7\ (12)}}
 \begin{cases}
 (p-1)/(p+1),&p\equiv1\pmod4,\\
 (p+1)/(p-1),&p\equiv3\pmod4
 \end{cases}
 =\frac{\sqrt3}{2}.
\]
\end{theorem}

Define the two real nonprincipal characters
\[
 \chi_4(n)=\begin{cases}0,&2\mid n,\\1,&n\equiv1\ (4),\\-1,&n\equiv3\ (4),\end{cases}
 \qquad
 \chi_{-3}(n)=\begin{cases}0,&3\mid n,\\1,&n\equiv1\ (3),\\-1,&n\equiv2\ (3).\end{cases}
\]
For $X\ge3$, put
\[
 L_X(1,\chi)=\prod_{p\le X}\left(1-\frac{\chi(p)}p\right)^{-1},
 \qquad
 P(X)=\prod_{\substack{p\le X\\p\equiv5,7\ (12)}}
 \left(\frac{p-1}{p+1}\right)^{\chi_4(p)}.
\]
We first work with finite products. At every prime $p>3$ the local
factor of $L_X(1,\chi_{-3})/L_X(1,\chi_4)$ is
\[
 \frac{1-\chi_4(p)/p}{1-\chi_{-3}(p)/p}.
\]
It is one for $p\equiv1,11\pmod{12}$, and is respectively
$(p-1)/(p+1)$ and its reciprocal for $p\equiv5,7\pmod{12}$.
The two exceptional factors at $p=2,3$ are $2/3$ and $4/3$.
Thus the exact identity is
\begin{equation}\label{eq:finite}
 P(X)=\frac98\frac{L_X(1,\chi_{-3})}{L_X(1,\chi_4)}
 \qquad(X\ge3).
\end{equation}

\section{Convergence and evaluation}
We justify the ordinary prime-ordered limit, rather than only an
Abelian limit with exponent $s\downarrow1$.
\begin{lemma}\label{lem:convergence}
For either character $\chi=\chi_4,\chi_{-3}$,
\[
 \lim_{X\to\infty}L_X(1,\chi)=L(1,\chi)
 =\sum_{n\ge1}\frac{\chi(n)}n.
\]
\end{lemma}
\begin{proof}
We use the established prime-number theorem estimate
\[
 \theta(x;q,a)=\frac{x}{\varphi(q)}+O_q(x/\log x),
 \qquad (a,q)=1,
\]
for $q=3,4$, supplied, for example, by Theorem 1.2 of
Bennett--Martin--O'Bryant--Rechnitzer~\cite{bmor}.
Since a nonprincipal character sums to zero over its unit classes,
\[
 \Theta_\chi(x):=\sum_{p\le x}\chi(p)\log p=O(x/\log x).
\]
Partial summation with weight $1/(t\log t)$ gives
\[
 \sum_{p\le X}\frac{\chi(p)}p
 =\frac{\Theta_\chi(X)}{X\log X}
 +\int_2^X\Theta_\chi(t)\frac{\log t+1}{t^2(\log t)^2}\,dt.
\]
The endpoint tends to zero and the integral converges, its large-$t$
integrand being $O(1/(t(\log t)^2))$. Hence the prime sum converges.
The higher terms in
\[
 -\log(1-\chi(p)/p)=\frac{\chi(p)}p+
       \sum_{k\ge2}\frac{\chi(p)^k}{kp^k}
\]
are absolutely summable over primes. The logarithms, and therefore
the positive partial Euler products, converge to a nonzero limit.
Abel's theorem applied to the convergent prime sum, with weights
$p^{-(s-1)}$, identifies the logarithmic limit as
$\lim_{s\downarrow1}\log L(s,\chi)$: the higher-power terms pass
to the limit by absolute domination. Finally, bounded partial sums
of the periodic nonprincipal character give continuity of its
Dirichlet series from the right at $s=1$. The Euler product for
$s>1$ therefore identifies the limit with $L(1,\chi)$.
\end{proof}

Both required values have elementary integral evaluations. Pairing
terms of the defining character series and summing their geometric
series under the integral gives
\begin{align*}
 L(1,\chi_4)&=\sum_{j\ge0}\left(\frac1{4j+1}-\frac1{4j+3}\right)
   =\int_0^1\frac{dt}{1+t^2}=\frac\pi4,\\
 L(1,\chi_{-3})&=\sum_{j\ge0}\left(\frac1{3j+1}-\frac1{3j+2}\right)
   =\int_0^1\frac{dt}{1+t+t^2}=\frac\pi{3\sqrt3}.
\end{align*}
The paired integrands are nonnegative, justifying these exchanges
by monotone convergence. Taking the limit in \eqref{eq:finite}
and using Lemma~\ref{lem:convergence} now proves Theorem~\ref{thm:main}:
\[
 \lim_{X\to\infty}P(X)=\frac98\frac{\pi/(3\sqrt3)}{\pi/4}
 =\frac{\sqrt3}{2}.
\]

The accompanying exact verifier checks all 1000 published prime
terms and the local quadratic-residue classification, the rational
identity \eqref{eq:finite} at every prime cutoff through $10000$,
and the exceptional factors separately. Numerical partial products
are supplementary controls; convergence follows from the analytic
argument and the cited prime-number theorem estimate.

\begin{thebibliography}{2}
\bibitem{oeis} The OEIS Foundation Inc.,
\href{https://oeis.org/A003630}{A003630}, accessed October 27, 2026.
\bibitem{bmor} M. A. Bennett, G. Martin, K. O'Bryant and A. Rechnitzer,
\emph{Explicit bounds for primes in arithmetic progressions},
Theorem 1.2, \href{https://arxiv.org/abs/1802.00085}{arXiv:1802.00085},
version 3, 2018.
\end{thebibliography}
\end{document}
